\chapter{Evaluation and Results}

\section{Data Provenance and Limitations}

Before presenting results, we must clearly state the evaluation data sources:

\subsection{Simulated Rig Results}

Results in Section~\ref{sec:rig_results} are generated from \textbf{simulated telemetry} (\texttt{sim/telemetry.py}). The simulation creates synthetic CloudWatch metrics and X-Ray traces that mimic AWS patterns but are not actual AWS data.

\textbf{Purpose:} Validate that the detection/localisation/evaluation pipeline executes correctly end-to-end.

\textbf{Validity:} These results demonstrate pipeline correctness but do \textit{not} represent live AWS performance. They verify that:
\begin{itemize}
    \item The detector correctly fires on simulated metric anomalies
    \item The ranker correctly ranks services based on simulated trace patterns
    \item The evaluator correctly computes P/R/F1/Top-k from ground truth
\end{itemize}

\textbf{Limitation:} Real AWS telemetry includes noise, cold-start variance, network jitter, and correlated anomalies that simulation cannot fully capture. Live AWS F1 would likely be 5-10 percentage points lower than simulated results.

\subsection{RCAEval Benchmark Results}

Results in Section~\ref{sec:rcaeval_results} are from the \textbf{public RCAEval dataset} \citep{pham2025rcaeval}, containing real telemetry from Online Boutique microservices with 90 labeled fault cases.

\textbf{Validity:} These are legitimate like-for-like comparisons on common benchmark data shared across the research community.

\subsection{Overhead Results}

Overhead results in Section~\ref{sec:overhead_results} are from the \textbf{live AWS} lite Leg~3 campaign (\texttt{results/live/overhead.json}). Earlier estimator-only figures are superseded by these measured values.

\section{Rig Results: Simulated Telemetry}
\label{sec:rig_results}

\textit{Source: Simulated telemetry - functional check, NOT AWS data. Results verify pipeline correctness.}

\subsection{Detection Performance}

Table~\ref{tab:detection_sim} presents detection metrics on 48 simulated injections (4 fault types $\times$ 2 load levels $\times$ 6 replications).

\begin{table}[h]
\centering
\caption{Detection Performance on Simulated Telemetry}
\label{tab:detection_sim}
\begin{tabular}{lcccccc}
\hline
\textbf{Load} & \textbf{Fault} & \textbf{Inj.} & \textbf{TP} & \textbf{FP} & \textbf{Precision} & \textbf{Recall} & \textbf{F1} \\
\hline
All & All & 48 & 48 & 3 & 0.941 & 1.000 & 0.970 \\
Peak & All & 24 & 24 & 1 & 0.960 & 1.000 & 0.980 \\
Steady & All & 24 & 24 & 2 & 0.923 & 1.000 & 0.960 \\
\hline
All & dependency\_failure & 12 & 12 & 0 & 1.000 & 1.000 & 1.000 \\
All & elevated\_latency & 12 & 12 & 0 & 1.000 & 1.000 & 1.000 \\
All & throttling & 12 & 12 & 0 & 1.000 & 1.000 & 1.000 \\
All & timeout & 12 & 12 & 0 & 1.000 & 1.000 & 1.000 \\
\hline
\end{tabular}
\end{table}

\textbf{Observations:}
\begin{itemize}
    \item Recall is perfect (1.0): all 48 injections were detected
    \item Precision is high (0.94): only 3 false positives occurred during control/recovery periods
    \item Per-fault-type F1 is 1.0 for all types in simulated data (overly optimistic)
\end{itemize}

\textbf{Detection delay:} Median 60s, P95 60s. This matches the 1-minute CloudWatch aggregation period—detections fire as soon as the first metric datapoint exceeds thresholds.

\subsection{Localisation Performance}

Table~\ref{tab:localisation_sim} presents localisation accuracy on the 48 detected cases.

\begin{table}[h]
\centering
\caption{Localisation Performance on Simulated Telemetry}
\label{tab:localisation_sim}
\begin{tabular}{lccccc}
\hline
\textbf{Load} & \textbf{Fault} & \textbf{Top-1} & \textbf{Top-3} & \textbf{Mean Rank} & \textbf{End-to-End Top-3} \\
\hline
All & All & 0.917 & 0.979 & 1.15 & 0.979 \\
Peak & All & 0.917 & 0.958 & 1.17 & 0.958 \\
Steady & All & 0.917 & 1.000 & 1.12 & 1.000 \\
\hline
All & dependency\_failure & 1.000 & 1.000 & 1.00 & 1.000 \\
All & elevated\_latency & 1.000 & 1.000 & 1.00 & 1.000 \\
All & throttling & 1.000 & 1.000 & 1.00 & 1.000 \\
All & timeout & 0.667 & 0.917 & 1.58 & 0.917 \\
\hline
\end{tabular}
\end{table}

\textbf{Observations:}
\begin{itemize}
    \item Top-3 accuracy is near-perfect (0.98) overall
    \item Timeout faults are harder to localise (Top-1 = 0.67): the target service times out but the symptom (504 error) appears at the caller, confusing the ranker
    \item Mean rank is very low (1.15): the ground-truth service typically ranks first or second
\end{itemize}

\textbf{Interpretation:} Simulated telemetry is clean (no noise, no correlated anomalies). Real AWS performance would likely show lower Top-3 (0.85-0.90) and higher mean rank (1.5-2.0).

\section{RCAEval Benchmark Results}
\label{sec:rcaeval_results}

\textit{Source: Public RCAEval RE2-OB raw under \texttt{results/rcaeval/}; aggregated by \texttt{eval/leg2.py} into \texttt{results/rcaeval/summary.md} / \texttt{localisation.csv} / \texttt{detection.json}. Offline benchmark --- not live AWS.}

\subsection{Baseline Method Comparison}

Table~\ref{tab:rcaeval_comparison} reports AC@1 / AC@3 / mean rank from committed raw. CausalRCA is incomplete ($n{=}4$); do not treat it as a full 90-case peer.

\begin{table}[h]
\centering
\caption{Localisation on RCAEval raw (authoritative; CausalRCA $n{=}4$)}
\label{tab:rcaeval_comparison}
\begin{tabular}{lcccc}
\hline
\textbf{Method} & \textbf{$n$} & \textbf{AC@1} & \textbf{AC@3} & \textbf{Mean rank} \\
\hline
\textbf{Rule-based (ours)} & 90 & 0.289 & 0.611 & 2.98 \\
BARO & 90 & 0.144 & 0.878 & 2.32 \\
CIRCA & 90 & 0.589 & 0.878 & 2.08 \\
TraceRCA & 90 & 0.111 & 0.644 & 2.92 \\
Hybrid (optional) & 90 & 0.256 & 0.478 & 3.63 \\
CausalRCA & \textbf{4} & 0.000 & 1.000 & 3.00 \\
\hline
\end{tabular}
\end{table}

\textbf{Key findings (evidence-locked):}
\begin{itemize}
    \item Rules AC@3 $=0.611$ vs BARO/CIRCA $=0.878$ --- gap $26.7$\,pp (exceeds the pre-committed $>$10\,pp concession bar)
    \item Hybrid AC@3 $=0.478$ is \textbf{worse} than rules; hybrid is beyond CA2 primary scope
    \item CausalRCA cells are quarantined to $n{=}4$ with \texttt{fixed\_order.json} share$=1.0$; ranking tables that implied $n{=}90$ for CausalRCA are withdrawn; CausalRCA is \textbf{not} a peer baseline and is not an AWS blocker
\end{itemize}

\subsection{Hypothesis Test: Detection Accuracy}

RCAEval baselines are localisers, not detectors. Rule-arm detection on committed raw ($n{=}90$): TP$=90$, FP$=204$, FN$=0$ $\Rightarrow$ precision $=0.306$, recall $=1.0$, \textbf{F1 $=0.469$} (\texttt{detection.json}).

\begin{itemize}
    \item Xing et al.\ (2025) citation ceiling F1 $=0.938$ is Leg~1 literature only
    \item Gap $0.938-0.469=46.9$\,pp --- \textbf{not} within 10\,pp
    \item The withdrawn ``F1$\sim$0.85 from Top-3'' shortcut is invalid
\end{itemize}

\subsection{Hypothesis Test: Localisation Accuracy}

\textbf{H2:} Rule Top-3 concedes $>$10\,pp vs strongest baselines.

Bootstrap 95\% CI for (CIRCA$-$rules) Top-3 on $n{=}90$ $= [0.156,\ 0.378]$ (\texttt{summary.md}) --- the $>$10\,pp concession expectation is \textbf{supported}.

\subsection{Per-Fault-Type Breakdown}

Table~\ref{tab:rcaeval_faults} uses AC@3 from \texttt{localisation\_by\_fault.csv} (rules vs BARO; 15 cases each).

\begin{table}[h]
\centering
\caption{AC@3 by fault type (RCAEval raw)}
\label{tab:rcaeval_faults}
\begin{tabular}{lccc}
\hline
\textbf{Fault Type} & \textbf{Cases} & \textbf{Rule-based} & \textbf{BARO} \\
\hline
CPU & 15 & 0.667 & 0.800 \\
Memory & 15 & 0.600 & 1.000 \\
Delay & 15 & 0.600 & 0.800 \\
Loss & 15 & 0.533 & 0.867 \\
Disk & 15 & 0.600 & 1.000 \\
Socket & 15 & 0.667 & 0.800 \\
\hline
\end{tabular}
\end{table}

\section{Overhead Analysis}
\label{sec:overhead_results}

\textit{Source: \textbf{Live AWS Leg~3 lite} (\texttt{results/live/overhead.json}; raw \texttt{data/runs/live\_lite\_overhead/}). Protocol: \texttt{configs/experiment\_lite\_overhead.yaml} --- 5\,min $\times$ 3 conditions at 1\,rps (300 requests/condition) on the \texttt{faultlab} stack in \texttt{eu-west-1}. Lite is directional measured evidence, not confirmatory 30-minute cells. Stack destroyed after the round.}

\subsection{Telemetry Volume (live lite)}

Table~\ref{tab:overhead_volume} reports CloudWatch Logs IncomingBytes plus X-Ray trace bytes per 1000 requests under the three telemetry conditions.

\begin{table}[h]
\centering
\caption{Measured telemetry volume per 1000 requests (lite Leg~3)}
\label{tab:overhead_volume}
\begin{tabular}{lrrr}
\hline
\textbf{Condition} & \textbf{bytes/1000} & \textbf{traces} & \textbf{log bytes} \\
\hline
Full (Active, INFO, sample 1.0) & $1.900\times10^{7}$ & 511 & $1.271\times10^{6}$ \\
Policy (Active, ERROR, sample 0.05) & $3.740\times10^{6}$ & 51 & $7.370\times10^{5}$ \\
Off (PassThrough, ERROR) & $1.252\times10^{6}$ & 0 & $3.755\times10^{5}$ \\
\hline
\textbf{reduction policy vs full} & \multicolumn{3}{r}{$0.803$ ($\geq 0.5$ expected --- supported)} \\
\hline
\end{tabular}
\end{table}

\subsection{Latency / cost (live lite)}

End-to-end load-generator latency (policy vs off): Mann--Whitney $p\approx1.2\times10^{-5}$, median difference $\approx -42$\,ms (policy lower). List-price cost per million requests from measured volumes: full $\approx\$10.93$, policy $\approx\$2.25$. Learned lower bound (compressed RCAEval parquet, $n{=}8$ RE2-OB \texttt{checkoutservice} cases): median $4.23\times10^{5}$ bytes/1000 requests --- a lower bound on learned-arm inputs, not a live-service measurement (\texttt{results/live/learned\_lower\_bound.csv}).

\section{Competitiveness Assessment}

Pre-committed rule: competitive if F1 within 10\,pp of the strongest baseline \textbf{and} telemetry volume $\geq$50\% reduced \textbf{on live overhead}.

\begin{itemize}
    \item Rule-arm F1 $=0.469$ vs Xing ceiling $0.938$ --- \textbf{fails} the 10\,pp bar
    \item Rules AC@3 $=0.611$ vs BARO/CIRCA $=0.878$ --- localisation also concedes $>$10\,pp
    \item Overhead $\geq$50\% reduction: \textbf{supported on lite} (reduction $=0.803$)
\end{itemize}

\textbf{Conclusion:} Live lite overhead meets the telemetry-reduction limb, but F1/AC@3 still fail the accuracy limb. The approach is therefore \textbf{not} competitive under the joint CA2 decision rule. Sim F1$\approx$0.97 must not be substituted for RCAEval.

\section{Implications for Research and Practice}

\subsection{For Researchers}
\begin{itemize}
    \item Three-leg protocol remains useful; Leg~2 raw already falsifies invented Top-k/F1 tables
    \item CausalRCA incompleteness ($n{=}4$) must be disclosed whenever that baseline is tabulated
    \item Hybrid IF arm underperformed rules on Top-3 --- negative result retained
\end{itemize}

\subsection{For Practitioners}
\begin{itemize}
    \item Artefact deployed for lite Leg~3 then destroyed; reuse \texttt{docs/LITE\_LEG3\_OVERHEAD.md}
    \item Budget from measured lite volumes: policy $\approx\$2.25$/M req vs full $\approx\$10.93$/M (list prices)
\end{itemize}

\section{Negative Results and Limitations}

\begin{itemize}
    \item Simulated rig Top-3 $0.979$ vs RCAEval rules $0.611$ --- sim overstates accuracy
    \item Lite Leg~3 only (5\,min cells); full 30-minute protocol not re-run
    \item CausalRCA $n{=}4$ only
    \item Live F1 under real CloudWatch noise is unknown; prior $0.88$--$0.92$ guesses are withdrawn
\end{itemize}

\section{Summary}

\textbf{RQ answer (evidence-locked):} On RCAEval raw, rule-arm F1 $=0.469$ and AC@3 $=0.611$ (vs BARO/CIRCA $0.878$). Lite Leg~3 live overhead shows policy vs full volume reduction $=0.803$ (\texttt{results/live/overhead.json}). The joint competitiveness rule still fails on accuracy; claim hygiene prefers this negative positioning over invented wins.
